\documentclass[../ICIT-Paper-2024.tex]{subfiles}
\usepackage{hyperref} 
\usepackage{tablefootnote}
\begin{document}
\section{CONCLUSION}\label{section5}
This study applied various optimization techniques including Stochastic Gradient Descent (SGD), Adam Optimization Algorithm, Genetic Algorithm (GA), and Multilayer Perceptron (MLP) to optimize the credit scoring formula for electronic wallets on the DeFi platform. Experimental results showed that the Multilayer Perceptron achieved the highest efficiency compared to the other techniques. However, since the main objective of the study is to determine the optimal parameter set for the credit scoring formula, the Genetic Algorithm is considered a more suitable choice. The reasons for this choice are as follows:

\begin{itemize} \item Genetic Algorithm provides higher interpretability compared to Multilayer Perceptron, making it easier for users to understand the basis for the scoring results. \item Genetic Algorithm can produce a specific score compared to MLP, which only provides score ranges. \item Genetic Algorithm can integrate conditions for the parameters, allowing conditions to be changed clearly.  \end{itemize}

The study also indicates that using public blockchain data can significantly benefit optimizing machine learning models. This data provides a rich source of information and easy access to build and evaluate models, enhancing the efficiency and transparency of decentralized financial systems. In the future, we will expand this data from None-EVM chains like Cosmos~\footnote{\href{https://cosmos.network/}{https://cosmos.network}}, Solana~\footnote{\href{https://solana.com/}{https://solana.com}}, Polkadot~\footnote{\href{https://polkadot.network/}{https://polkadot.network}}, etc. We also focus on establishing an efficient mechanism for labeling this data. Based on labeled data, we do a backtest for optimizing models to demonstrate that our score can be applied in real life.
\end{document}